\chapter{Command Cheat Sheet}\label{app:cheatsheet}

The commands you will type most often, grouped by phase. Tool options change
between versions; when one fails, check the tool's own documentation for the
version you installed (and record that version in your notebook).

\section*{Python environment and metrics (Phase 0)}
\begin{lstlisting}[style=shell]
python3 -m venv .venv && source .venv/bin/activate
pip install -e ../Intention_based_speech_enhancment_using_NLP   # thesis code
pip install numpy scipy soundfile pandas matplotlib pytest ruff pesq pystoi pyyaml cocotb
pip freeze > requirements-lock.txt
pytest metrics                       # metric unit tests
ruff check .                         # lint Python
PYTHONPATH=src python docs/research_guide/scripts/measure_workload.py   # (thesis repo)
\end{lstlisting}

\section*{Golden model and sweeps (Phase 1)}
\begin{lstlisting}[style=shell]
python -m golden.sweep --config experiments/012-wordlength/config.yaml
python -m golden.vectors --out golden/vectors/      # export golden test vectors
git tag spec-v1 && git push --tags
\end{lstlisting}

\section*{RTL: lint, simulate, synthesize (Phases 2--3)}
\begin{lstlisting}[style=shell]
verilator --lint-only -Wall rtl/*.sv --top-module ise_top
make -C tb/recip_unit SIM=verilator          # one cocotb testbench (Makefile flow)
make -C tb/recip_unit SIM=icarus             # same test, second simulator
pytest tb/                                   # if you use cocotb's Python runner
yosys -p "read_verilog -sv rtl/bin_pipeline.sv; synth -top bin_pipeline; stat"
sv2v rtl/*.sv > build/design.v               # if a tool needs plain Verilog
sby -f tb/formal/min_tracker.sby             # formal proof (SymbiYosys)
gtkwave build/dump.fst  # or: surfer build/dump.fst
\end{lstlisting}

\section*{Coverage (Phase 3)}
\begin{lstlisting}[style=shell]
# build the simulation with Verilator's --coverage, run the tests, then:
verilator_coverage --annotate build/cov_annotated coverage.dat
\end{lstlisting}

\section*{Vivado in batch mode (Phase 4)}
\begin{lstlisting}[style=shell]
vivado -mode batch -source fpga/pynq_z2/build.tcl      # rebuild project + bitstream
\end{lstlisting}
\begin{lstlisting}[style=tcl]
# useful commands inside a Vivado Tcl script or console
report_timing_summary -file reports/timing.rpt
report_utilization    -file reports/utilization.rpt
report_power          -file reports/power.rpt
write_bd_tcl -force fpga/pynq_z2/bd.tcl                 # save the block design
\end{lstlisting}

\section*{On the PYNQ board (Phase 4)}
\begin{lstlisting}[style=py]
from pynq import Overlay, MMIO, allocate
ol = Overlay("ise.bit")              # .hwh with the same name next to it
print(ol.ip_dict.keys())             # find your IP and its address
regs = MMIO(ol.ip_dict["regs_axil_0"]["phys_addr"], 0x100)
print(hex(regs.read(0x08)))          # VERSION register: a first sanity check
\end{lstlisting}

\section*{ASIC flow (Track B)}
\begin{lstlisting}[style=shell]
nix-shell librelane/shell.nix        # enter the LibreLane environment
librelane --smoke-test
librelane asic/bin_pipeline/config.json            # default PDK
librelane --pdk ihp-sg13g2 asic/bin_pipeline/config.yaml   # another PDK
klayout runs/<run>/final/gds/bin_pipeline.gds     # look at your layout
\end{lstlisting}

\section*{Git habits}
\begin{lstlisting}[style=shell]
git switch -c exp/023-log-trackers          # one branch per experiment
git commit -m "Log-domain trackers: 8-bit minima meet RQ1 on dev set"
git tag m2-bit-exact && git push --tags    # one tag per milestone
\end{lstlisting}

\section*{This guide}
\begin{lstlisting}[style=shell]
cd docs/research_guide && make              # build research_guide.pdf
tools/check_chapter.sh chapters/05-codesign # check one chapter quickly
\end{lstlisting}
